\begin{textsample}{POS Dim 2 – to – Score 24.00 – to/to\_em\_36.txt}  \label{ex:f2_pos_068}
2.2 A trajetória da Educação Profissional na Rede Estadual do Tocantins O \textbf{Governo} do Estado, através da Secretaria da Educação, sempre buscou oferecer e garantir ao cidadão tocantinense o acesso a uma Educação Profissional de qualidade, integrada à comunidade e ao mundo do trabalho.
A partir desse compromisso, deu -se início às discussões e ao planejamento para a implantação e expansão da Educação Profissional \textbf{Técnica} de Nível Médio nas unidades escolares da rede pública estadual.
As principais atividades econômicas produtivas do Estado são a pecuária bovina de corte, o plantio de arroz, abacaxi e soja.
A atividade industrial ainda é pequena, mas encontra -se em franco crescimento.
No setor de serviços, além do comércio, predominam o sistema financeiro, as atividades turísticas, os serviços e investimentos realizados pelo setor público, que ainda é o maior empregador da economia tocantinense.
O comércio se caracteriza por ser distribuidor de mercadorias oriundas de outros estados.
O \textbf{Governo} do Estado do Tocantins, através da Secretaria da Educação, assinou junto ao Ministério da Educação o Termo de Adesão para implantação do Ensino Médio Integrado ao Ensino Médio, dando início, em 2006, ao Projeto Piloto em dez unidades escolares, localizadas em sete \textbf{municípios} do Estado do Tocantins, com demanda para atendimento de 480 alunos.
A iniciativa buscava \textbf{atender} às escolas públicas da rede estadual, com base na Lei nº 9.394/96, de \textbf{Diretrizes} e Bases da Educação Nacional, que trata do tema da “ Educação Profissional ” em um \textbf{capítulo} especial, o \textbf{Capítulo} III do Título V, “ Da Educação Profissional ”, \textbf{Artigos} 39 e 42, de forma associada e articulada com o § 2º do \textbf{Artigo} 36 da mesma Lei, na parte referente ao Ensino Médio ; e no Decreto nº 5.154/2004, que, por seu turno, define que “ a Educação Profissional \textbf{Técnica} de nível médio e o Ensino Médio ” dar -se-ão de forma integrada, concomitante e subsequente ao Ensino Médio ( cf. incisos I, II e III do § 1º do \textbf{Artigo} 4º ).
Na ocasião, foram \textbf{ofertados} onze \textbf{cursos} \textbf{técnicos}, selecionados conforme o potencial de mercado da região.
O processo de implantação foi iniciado em dez escolas-piloto : Escola Estadual Agrícola David Aires França, em Arraias ; Estadual \textbf{Técnica} de Enfermagem de Araguaína, em Araguaína ; Centro de Ensino Médio Ary Ribeiro Valadão Filho, em Gurupi ; Escola Estadual Brigadas Che Guevara, em Monte do Carmo ; Escola \textbf{Técnica} Agrícola de Natividade, em Natividade do Tocantins ; Centro de Ensino Médio Florêncio Aires, em Porto Nacional ; Escola Estadual Dr.
José de Sousa Porto, em Pedro Afonso ; Centro de Ensino Médio Deputado Darcy Marinho, em Tocantinópolis ; e duas outras escolas-piloto onde a \textbf{oferta} dos \textbf{cursos} era realizada a partir de parcerias : a Escola \textbf{Técnica} de Saúde do SUS – ETSUS e a Escola \textbf{Técnica} Federal de Palmas, ambas em Palmas.
Entretanto, esta modalidade teve seu \textbf{fortalecimento} no Estado através da celebração do convênio com o Programa Brasil Profissionalizado, criado através do Decreto nº 6.302/2007.
Este programa foi \textbf{instituído} com a perspectiva de expandir e ampliar a \textbf{oferta} de \textbf{cursos} \textbf{técnicos} de nível médio, principalmente do ensino médio integrado à educação profissional e tecnológica, viabilizando a construção, \textbf{reforma} e modernização de unidades escolares, a aquisição de equipamentos, mobiliários e laboratórios, além do financiamento de recursos pedagógicos e de formação e \textbf{qualificação} dos profissionais da educação.
Nessa trajetória de implantação de \textbf{cursos} \textbf{técnicos} integrados ao ensino médio na rede estadual, três \textbf{instituições} \textbf{técnicas} estaduais ligadas à educação profissional foram transferidas para a pasta da então Secretaria de Ciência e Tecnologia do Tocantins em 2008.
Foram elas : o Colégio Estadual Agrícola de Pedro Afonso, a Escola \textbf{Técnica} Agrícola de Natividade e a Escola \textbf{Técnica} de Enfermagem de Araguaína.
No mesmo ano, a Escola \textbf{Técnica} Federal de Palmas herdou o \textbf{curso} \textbf{Técnico} em Enfermagem da então Estadual \textbf{Técnica} de Enfermagem de Araguaína, por meio de um Termo de Cooperação \textbf{Técnica} entre a Secretaria de Ciência e Tecnologia do Estado do Tocantins e a antiga Escola \textbf{Técnica} Federal de Palmas.
Em junho de 2014, o \textbf{Governo} do Estado do Tocantins cedeu a área e as instalações do Colégio Estadual Agrícola Dr.
José de Souza Porto ao Instituto Federal do Tocantins, através do Decreto nº 5.037, de 9 de maio de 2014, passando a se chamar Campus Avançado Pedro Afonso.
Já a Escola \textbf{Técnica} Agrícola de Natividade retornou à pasta da Secretaria da Educação em 2015.
Com base no Censo Escolar de 2021, a rede estadual de ensino do Tocantins realizou atendimento a um número de 1.871 estudantes, em vinte unidades escolares com a \textbf{oferta} de \textbf{cursos} \textbf{técnicos} integrados ao ensino médio, na forma articulada ( desenvolvida nas formas integrada e concomitante ) e subsequente.
A Seduc vem realizando estudos para a ampliação gradativa da \textbf{oferta} da modalidade na rede estadual, onde são considerados os seguintes fatores : a escolha dos estudantes ; as demandas socioeconômicas e ambientais dos cidadãos e do mundo do trabalho ; o \textbf{fortalecimento} das parcerias ; e a capacidade de investimento para a expansão, reestruturação e manutenção dos \textbf{cursos}.
Os passos que permitirão a expansão da educação profissional no Estado serão dados com cautela e muito planejamento, para que essa modalidade possa cumprir a finalidade \textbf{prevista} na Lei de \textbf{Diretrizes} e Bases da Educação Nacional ( LDB ) de preparar “ para o exercício de profissões ”, contribuindo para que o cidadão possa se inserir e atuar no mundo do trabalho e na vida em sociedade.
Tabela.
Total de \textbf{cursos} do EMI, incluindo PROEJA, \textbf{ofertados} pela rede estadual de ensino ( 2021 ) A \textbf{oferta} da Educação Profissional \textbf{Técnica} de Nível Médio no território do Tocantins vem crescendo de forma gradativa.
Com base no Censo Escolar 2021, o estado do Tocantins atingiu 7.762 matrículas em \textbf{cursos} da Educação Profissional.
A tabela a seguir representa a divisão do total de matrículas entre as \textbf{ofertas} estadual, \textbf{federal}, \textbf{municipal} e privada.
A previsão de \textbf{oferta} de cinco diferentes \textbf{itinerários} a serem escolhidos pelos jovens, incluindo, entre eles, uma opção de formação \textbf{técnica} e profissional, representa uma das principais mudanças trazidas pela Lei nº 13.415/2017 e aparece como uma real possibilidade de \textbf{fortalecimento} da educação profissional.
A Lei citada propõe integrar a formação geral básica à formação profissional no \textbf{currículo} do Ensino Médio, \textbf{atendendo} às demandas ocasionadas pelas mudanças tecnológicas e pela organização do trabalho, que pedem profissionais com múltiplas habilidades e maior flexibilidade.
Espera -se que essa nova proposição represente um marco para a ampliação e o \textbf{fortalecimento} da formação \textbf{técnica} e profissional no território do Tocantins.

% matched lemmas: artigo, atender, capítulo, currículo, curso, diretriz, federal, fortalecimento, governo, instituir, instituição, itinerário, municipal, município, oferta, ofertar, prever, qualificação, reforma, técnico
\end{textsample}
